% Draft Section 5.2 prose and figure captions. Copy/adapt into the paper.
% Requires graphicx. Set the graphics path to this directory's figures/,
% or copy the four PDFs into the paper's figure directory.
% All empirical claims below refer to body8, floor, behavioral signals.

\paragraph{Transparency and cooperation under behavioral signals.}
We evaluate the eight-program library on the floor side at
$t\in\{0,0.01,\ldots,1\}$ and
$\alpha\in\{0,0.2,0.4,0.5,0.6,0.8,1\}$, with payoffs
$(T,R,P,S)=(3,2,0,-1)$. For each fixed positive caution level, maximum Nash
cooperation and the cooperation selected by the sampled replicator trajectories
are nondecreasing in transparency. Both change from zero to one at the same
sampled onset. For $\alpha=0.2,0.4,0.5,0.6,0.8,1$, respectively, the first
cooperative grid points are $t=0.28,0.01,0.11,0.24,0.50,1$.
At every cooperative point, all 1,000 sampled initial populations converge to
full mutual cooperation; at every remaining point with $\alpha>0$, all converge
to mutual defection. The maximum over symmetric equilibria equals the maximum
over all equilibria throughout the grid. Full-defection equilibria nevertheless
remain available at every positive caution level. Transparency therefore
enables cooperative equilibria and favors their selection in these experiments,
without eliminating defecting equilibria. At $\alpha=0$, cooperation is
unconditional at every transparency level. The 707 parameter points represent
62 distinct games, with the same 1,000 Dirichlet$(1)$ initial populations reused
in each game. The observed cooperative basin of one has a pointwise 95\%
Wilson interval $[0.9962,1]$; it is not a proof of global convergence.

\paragraph{Caution and selection.}
The transparency needed for cooperation depends nonmonotonically on caution.
Among the positive caution settings tested, $\alpha=0.4$ has the lowest sampled
onset, $t=0.01$, corresponding to approximately $0.03$ bits of information about
the opponent. The previous sampled point, $t=0$, has only defecting equilibria.
In contrast, $\alpha=1$ requires the full-transparency endpoint. Moreover, the
equilibrium effect cannot be read from the raw rate of cooperation before
selection. At $\alpha=0.2$, average initial cooperation falls from $57.65\%$ at
$t=0$ to $32.84\%$ at $t=1$, while cooperation after selection rises from zero
to one. All these comparisons are restricted to the sampled grid and the fixed
program library; the plotted transition brackets do not locate exact
continuous thresholds.

\begin{figure}[t]
  \centering
  \includegraphics[width=\linewidth]{01_transparency_and_cooperation.pdf}
  \caption{Transparency enables and selects mutual cooperation in the
  behavioral family. Left: maximum $\Pr[(C,C)]$ over the complete Nash set of
  each induced game, obtained by exact enumeration. Right: observed fraction
  of 1,000 uniformly sampled initial populations converging to full cooperation
  under replicator dynamics. Each stroke represents one sampled transparency
  value at the indicated caution level; circles mark the first cooperative
  grid point. The two quantities agree at all 707 parameter points. The same
  1,000 starts are reused across 62 distinct games. The pointwise 95\% Wilson
  intervals for basins observed as one and zero are $[0.9962,1]$ and
  $[0,0.0038]$, respectively. No interpolation across caution settings is
  implied. At $\alpha=0$, every program cooperates unconditionally.}
  \label{fig:behavioral-transparency-cooperation}
\end{figure}

\begin{figure}[t]
  \centering
  \includegraphics[width=\linewidth]{02_equilibrium_landscape.pdf}
  \caption{Nash possibilities and evolutionary selection at every tested
  positive caution level. Teal shows the Nash cooperation level reached by all
  sampled replicator runs; purple shows the other Nash cooperation levels.
  Every connected Nash component in these games has constant cooperation, so
  their union gives the complete Nash outcome set, including continua.
  Mutual-defection equilibria persist even when the dynamics select full
  cooperation. Shaded strips locate the interval between the last sampled
  defecting point and the first sampled cooperative point; the grid resolution
  is $0.01$. Lines guide the eye between observed values. At $\alpha=1$, the
  cooperative endpoint is shown separately.}
  \label{fig:behavioral-nash-landscape}
\end{figure}

\begin{figure}[t]
  \centering
  \includegraphics[width=\linewidth]{03_transparency_thresholds.pdf}
  \caption{The first cooperative transparency level is nonmonotone in caution.
  Dots mark the first sampled $t$ with a fully cooperative Nash equilibrium
  and 1,000/1,000 cooperative replicator outcomes; black ticks indicate the
  preceding defecting grid point. The right axis converts the nominal
  transparency to mutual information, $I(B;S)=t\log_2 8=3t$ bits, under the
  uniform calibration prior. Among the tested positive caution values,
  $\alpha=0.4$ requires the least sampled transparency, $t=0.01$, approximately
  $0.03$ bits. Dots do not identify exact continuous thresholds, and no
  interpolation between caution settings is assumed.}
  \label{fig:behavioral-transparency-onset}
\end{figure}

\begin{figure}[t]
  \centering
  \includegraphics[width=\linewidth]{04_cooperation_before_and_after_selection.pdf}
  \caption{Cooperation before and after evolutionary selection. For each
  game, gray shows the mean and interquartile range of $x^\top Cx$ over the
  1,000 Dirichlet$(1)$ starting populations, where $C_{ij}$ indicates a
  mutual-cooperation cell. Teal shows the mean at the converged endpoints,
  which is either zero or one in every game. Gray bands describe variation
  among starting populations, not confidence intervals. In particular,
  cooperation before selection can fall with transparency even while
  selected cooperation rises. The same starts are used in all games; lines
  connect sampled values and the $\alpha=1$ endpoint is shown separately.}
  \label{fig:behavioral-selection-amplification}
\end{figure}
